% ============================================================
% f2014_liaoning_science_q07_three_views.tex
% 2014 辽宁卷（理科）第 7 题：三视图
% 来源 PDF：普通高考/2014/2014辽宁理.pdf
% 原页：第 1 页（A3），印面 895，第 7 题题干下方
% 主视图、左视图均为 2×2 矩形，中部有竖线；
% 俯视图为 2×2 正方形去掉左上、右下两个半径 1 的四分之一圆。
% ============================================================
\begin{tikzpicture}[>=Stealth, line width=0.7pt, font=\small,
  dim/.style={line width=0.4pt},
  every node/.style={inner sep=1.5pt}]
  % ---------- 主视图 ----------
  \draw (0,3.0) rectangle (2,5.0);
  \draw (1,3.0) -- (1,5.0);
  \node at (1,2.62) {主视图};
  \draw[dim,<->] (0,5.2) -- (1,5.2) node[midway,above,font=\scriptsize] {$1$};
  \draw[dim,<->] (0,5.6) -- (2,5.6) node[midway,above,font=\scriptsize] {$2$};
  \draw[dim,<->] (-0.3,5.0) -- (-0.3,3.0) node[midway,left,font=\scriptsize] {$2$};

  % ---------- 左视图 ----------
  \draw (3.6,3.0) rectangle (5.6,5.0);
  \draw (4.6,3.0) -- (4.6,5.0);
  \node at (4.6,2.62) {左视图};
  \draw[dim,<->] (3.6,5.2) -- (4.6,5.2) node[midway,above,font=\scriptsize] {$1$};
  \draw[dim,<->] (3.6,5.6) -- (5.6,5.6) node[midway,above,font=\scriptsize] {$2$};

  % ---------- 俯视图 ----------
  % 正方形去掉左上、右下两个四分之一圆（半径 1）
  \draw (0,0) -- (1,0)
        arc[start angle=180, end angle=90, radius=1]
        -- (2,2) -- (1,2)
        arc[start angle=0, end angle=-90, radius=1]
        -- cycle;
  \node at (1,-0.95) {俯视图};
  \draw[dim,<->] (0,-0.32) -- (2,-0.32) node[midway,below,font=\scriptsize] {$2$};
  \draw[dim,<->] (2.35,0) -- (2.35,2) node[midway,right,font=\scriptsize] {$2$};
  % 两个半径标注
  \draw[dim,->] (0,2) -- (0.707,1.293) node[midway,above left,font=\scriptsize] {$1$};
  \draw[dim,->] (2,0) -- (1.293,0.707) node[midway,below right,font=\scriptsize] {$1$};
\end{tikzpicture}
